% Generated by mcm-paper-writer.
% This is a stable training/draft layout. Use the current official contest template for submission.
\documentclass[UTF8,12pt,a4paper,fontset=none]{ctexart}
\IfFontExistsTF{Songti SC}{
  \setCJKmainfont{Songti SC}[BoldFont={Songti SC Bold}]
  \setCJKsansfont{PingFang SC}[BoldFont={PingFang SC Semibold}]
  \setCJKmonofont{PingFang SC}
}{
  \IfFontExistsTF{SimSun}{
    \setCJKmainfont{SimSun}[BoldFont={SimHei}]
    \setCJKsansfont{Microsoft YaHei}[BoldFont={Microsoft YaHei Bold}]
    \setCJKmonofont{Microsoft YaHei}
  }{
    \IfFontExistsTF{Microsoft YaHei}{
      \setCJKmainfont{Microsoft YaHei}[BoldFont={Microsoft YaHei Bold}]
      \setCJKsansfont{Microsoft YaHei}[BoldFont={Microsoft YaHei Bold}]
      \setCJKmonofont{Microsoft YaHei}
    }{
      \IfFontExistsTF{Noto Serif CJK SC}{
        \setCJKmainfont{Noto Serif CJK SC}
        \setCJKsansfont{Noto Sans CJK SC}
        \setCJKmonofont{Noto Sans Mono CJK SC}
      }{
        \IfFontExistsTF{Source Han Serif SC}{
          \setCJKmainfont{Source Han Serif SC}
          \setCJKsansfont{Source Han Sans SC}
          \setCJKmonofont{Source Han Sans SC}
        }{
          \PackageError{mcm-paper-writer}
            {No renderer-compatible CJK font found}
            {Install Songti SC, SimSun, Microsoft YaHei, Noto CJK, or Source Han CJK before compiling}
        }
      }
    }
  }
}
\usepackage{geometry}
\geometry{top=2.5cm,bottom=2.5cm,left=2.5cm,right=2.5cm}
\usepackage{amsmath,amssymb,bm}
\usepackage{booktabs,array,tabularx}
\usepackage{graphicx,float}
\usepackage{caption}
\usepackage{enumitem}
\usepackage{setspace}
\usepackage[hidelinks]{hyperref}
\setstretch{1.2}
\setlength{\parindent}{2em}
\setlength{\parskip}{0.25em}
\renewcommand{\arraystretch}{1.25}
\pagestyle{plain}
\renewenvironment{abstract}{%
  \par\small
  \begin{center}\bfseries\abstractname\end{center}%
  \par
  \setlength{\leftskip}{0pt}%
  \setlength{\rightskip}{0pt}%
}{%
  \par
}
\title{无线电干扰源的快速自动定位与清除}
\author{}
\date{}

\begin{document}
\maketitle
\begin{center}题号：B\end{center}

\begin{abstract}
针对圆形目标区域内未知个数无线电干扰源的自动定位与清除，建立有界误差示向度的前向角域交会模型，并给出第二检测点选择、全向在线搜索与定向混合覆盖清除算法。问题1证明：以定位区域直径为直径的圆一般不能覆盖该区域；等边三角形反例的覆盖判定为否。问题2以最坏剩余覆盖半径评价第二检测点，理论扇形半径为750.114米。问题3主线GREEDY\_FAST在10局官方演练中全部清除，平均定位清除时间341秒，并用全清率检验策略有效性。问题4采用HEX37三角晶格覆盖，离线30组配对全部全清并完成无源证书，平均T/K为836.91秒；演练PREFIX\_A局T/K为822.38秒。上述数值均为演练或离线几何结果，正式测试尚未执行。

\noindent\textbf{关键词：}交会定位；示向度；覆盖路径；在线搜索；定向干扰源
\end{abstract}

\clearpage

\section{问题背景与问题重述}

无线电干扰源的精准定位通常要在已知分布区域内近距测向完成。本题给出半径1800米的圆形目标区域，源个数在10到16之间但在线未知，要求一只机器狗从原点出发，利用测向、移动与光学清除接口，在全部清除的前提下缩短虚拟任务时间。

测向机在覆盖内返回示向度，误差不超过1度；过近则只报告near，无信号则报告no\_signal。清除成功当且仅当指定频道存在未清除源且欧氏距离不超过20米，与射频朝向无关。问题4另有定向源，有效辐射为朝向两侧各90度的闭扇形。

四个子问题依次为：由多点示向度计算交会多边形直径并回答直径圆能否覆盖；在仅有一个全向示向度时给出第二检测点策略与候选区；对全部全向源给出在线搜索定位清除算法并用演练检验；在全向与定向混合且朝向未知时扩展策略。附件给出接口计时与两局样例会话，清洗后得到12条有效动作，其中direction 2条、no\_signal 4条，用于核对观测形态，不作为成绩。

本文不把1800米圆当作机器狗禁出区，不把清满10个当作结束，不把示向度当精确射线，也不把单点全频道无信号当作无源证书。覆盖路径规划的一般方法见文献，本题的证书必须针对未知朝向与未知接收半径专门证明。

\section{模型假设与符号}

假设干扰源位置在目标圆内固定，频道互异且时不变；全向源覆盖为圆盘，定向源覆盖为180度闭扇形。示向误差为有界集合而非高斯分布，因此定位对象是可行集而不是点估计置信圆。机器狗允许离开目标圆，源不得在圆外。

\begin{table}[!htbp]
  \centering
  \caption{主要符号}
  \label{tab-symbols}
  \begin{tabularx}{\textwidth}{>{\centering\arraybackslash}X>{\centering\arraybackslash}X>{\centering\arraybackslash}X}
    \toprule
    符号 & 含义 & 单位 \\
    \midrule
    Ω & 目标圆域，半径1800 & m \\
N & 干扰源总数，属于10到16 & 1 \\
K & 已成功清除数 & 1 \\
θ̂ & 示向度，误差界1度 & deg \\
r\_i & 源i接收半径，属于[1000,1500] & m \\
ρ & 可行集到候选点的豪斯多夫半径 & m \\
T & 虚拟任务时间 & s \\
    \bottomrule
  \end{tabularx}
\end{table}

计时按附件协议：移动速度5米每秒，检测5秒，频道切换1秒，清除失败3秒，成功再加2秒。T包含空扫与失败动作。平均定位清除时间为T/K，K=0时无定义。

\section{数据来源与接口观测}

数值常数全部来自冻结参数表，共28项，与题面及附件一致。目标圆半径1800米、接收半径1000至1500米、示向误差1度、近距5米、清除20米。代码禁止另写一套常数。

附件样例会话经严格JSON解析后得到2局12条清洁记录，重复动作按幂等回放识别，未静默删行。这些记录只用于确认direction、near、no\_signal三种返回及计时字段，不能外推正式成绩。

问题3、问题4的数量结果来自官方演练接口与本地几何模拟器。演练各局随机非配对；离线30组使用固定种子与有界确定误差。正式测试日志不存在，表1相应单元格标记为无正式日志。

\section{模型建立}

问题1的定位区域是各检测点前向闭角域的交集。对检测点s与示向度α，误差界ε=1°，角域由两条过s的射线围成，等价为两个闭半平面。

\begin{equation}
W(s,\alpha)=\{\,x\in\mathbb{R}^{2}: n_{-}(x-s)\le 0,\; n_{+}(x-s)\le 0\,\}
\label{eq-wedge}
\end{equation}

多个检测点时取交集R。R可能为空、无界、退化为点或线段，算法必须输出显式状态。直径定义为区域内最远两点距离，对凸多边形只需检查顶点对。

\begin{equation}
\mathrm{diam}(R)=\max_{p,q\in R}\|p-q\|
\label{eq-diam}
\end{equation}

以直径为直径的圆一般不能覆盖原区域。等边三角形直径等于边长d，最小覆盖圆半径为d/√3>d/2，因此直径圆不能覆盖。该反例同时说明：清除点充分条件应使用到区域的豪斯多夫半径不超过20米，不能用直径的一半代替。

问题2只有一个全向示向度。第一点观测后，源位于楔形、目标圆与最大接收圆盘的交集F1。理论扇形可用一个半径

\begin{equation}
\rho_{\mathrm{sec}}=\frac{1500}{2\cos 1^{\circ}}\approx 750.114\,\mathrm{m}
\label{eq-sector-rho}
\end{equation}

的圆覆盖。第二点q的评价不是未知真位置上的最优布置，而是对near、direction与no\_signal三分支取最坏剩余覆盖半径J(q)。候选在F1最小覆盖圆附近分层生成，删除20至3000米固定极坐标网格。输出是候选集推荐，不是声称全局最优的q*。

问题3全部为全向源。发现层需要有限测点保证：对任意源位与未知半径下界1000米，至少一点落入接收圆。实现采用原点加半径1200米六等分外围点的七点集作为历史对照，主线在线策略为GREEDY\_FAST：对已发现频道优先用累计楔形收缩后的确定清除，失败则有限启发式，再回退光学网格SAFE。终止不能用K=10，必须对未观测频道给出覆盖证书或N=16计数证书。

问题4增加未知朝向。SQUARE81是600米方格81点，可证对任意g∈Ω与任意朝向d，存在格点落入最小接收半径内部且严格位于180度发射半平面内。HEX37改用边长800米三角晶格、环数3共37点，无限晶格覆盖半径

\begin{equation}
\delta=800/\sqrt{3}\approx 461.880\,\mathrm{m}
\label{eq-hex-delta}
\end{equation}

构造点y=g+500d落在半径2300米盘内，其最近格点同时满足距离与半平面松弛。有限37点相对第4环的归属有约9.4米间隙。同一点集覆盖全向源。频道在全部覆盖点无信号则判空。寻路使用离线Hamilton路，禁止在线k!枚举。

覆盖选址思想与最大覆盖问题相关，但本题证书针对最坏几何而非需求权重。最短路只用于覆盖图上的固定骨架，不把动态网络最短路当作最优搜索。

\section{算法与求解结果}

问题1算法：构造半平面，线性规划判定可行与无界；枚举相邻约束交点得到顶点，计算最远点对作为直径，并比较最小覆盖圆半径与直径一半。离线演示中等边三角形直径约为1，最小覆盖圆半径约为0.577，直径圆覆盖判定为否。

\begin{table}[!htbp]
  \centering
  \caption{问题1离线几何演示}
  \label{tab-q1-demo}
  \begin{tabularx}{\textwidth}{>{\centering\arraybackslash}X>{\centering\arraybackslash}X>{\centering\arraybackslash}X>{\centering\arraybackslash}X>{\centering\arraybackslash}X}
    \toprule
    实例 & 状态 & 直径 & 最小覆盖圆半径 & 直径圆是否覆盖 \\
    \midrule
    等边三角形 & polygon & 1.000 & 0.577 & 否 \\
两楔交会 & polygon & 9.877 & 4.939 & 演示未作覆盖断言 \\
    \bottomrule
  \end{tabularx}
\end{table}

\begin{figure}[!htbp]
  \centering
  \includegraphics[width=0.82\textwidth]{\detokenize{figures/fig-q1-diameter.png}}
  \caption{等边三角形定位区域：实线直径圆不能覆盖顶点，虚线为直径圆，深色圆为最小覆盖圆。}
  \label{fig-q1-diameter}
\end{figure}

问题2算法在F1外包多边形上计算最小覆盖圆c1，于半径比例0、0.25、0.5、0.75、1及探索环带上取候选，按J(q)排序。演示实例理论扇形半径750.114米，推荐点对应J约为43.03米，该值只说明候选集分辨率，不是现场第二点。

问题3在线流程：按发现点扫描频道；direction进入累计楔形后验；rho不超过20米则确定清除；否则有限中心试清与补测；失败进入光学SAFE。GREEDY\_FAST令SourceBudgetS=0，不在发现完成后放弃未清源。对照GREEDY全清但更慢，GREEDY\_ABORT放弃SAFE后全清仅1/10，不作候选。

\begin{table}[!htbp]
  \centering
  \caption{问题3官方演练10局对照（非正式、非配对）}
  \label{tab-q3-practice}
  \begin{tabularx}{\textwidth}{>{\centering\arraybackslash}X>{\centering\arraybackslash}X>{\centering\arraybackslash}X>{\centering\arraybackslash}X>{\centering\arraybackslash}X}
    \toprule
    策略 & 全清 & 平均T/K & 区间 & SAFE均值 \\
    \midrule
    GREEDY & 10/10 & 390 & 334-455 & 1.20 \\
GREEDY\_FAST & 10/10 & 341 & 244-450 & 0.20 \\
GREEDY\_ABORT & 1/10 & 382 & 295-482 & 0 \\
    \bottomrule
  \end{tabularx}
\end{table}

\begin{figure}[!htbp]
  \centering
  \includegraphics[width=0.82\textwidth]{\detokenize{figures/fig-q3-practice-tk.png}}
  \caption{问题3演练平均T/K。GREEDY\_FAST为341秒，柱面斜线、网格与点纹区分策略，不单靠颜色。}
  \label{fig-q3-practice-tk}
\end{figure}

问题4定位累计全部direction前向楔、1500米外包、目标圆外包与near的5米约束；Q4的no\_signal不挖距离圆盘。包围半径不超过20米时确定清除；光学兜底将凸定位域划为边长28米方格，半对角线14√2<20米。发现层默认HEX37，SQUARE81保留为回归基线。

\begin{table}[!htbp]
  \centering
  \caption{问题4离线配对种子0-29}
  \label{tab-q4-offline}
  \begin{tabularx}{\textwidth}{>{\centering\arraybackslash}X>{\centering\arraybackslash}X>{\centering\arraybackslash}X>{\centering\arraybackslash}X>{\centering\arraybackslash}X}
    \toprule
    覆盖 & 全清且证书 & 平均T/K (s) & 平均路程 (m) & 平均测向次数 \\
    \midrule
    SQUARE81 & 30/30 & 1313.16 & 68250.75 & 679.73 \\
HEX37 & 30/30 & 836.91 & 46196.45 & 337.10 \\
    \bottomrule
  \end{tabularx}
\end{table}

\begin{figure}[!htbp]
  \centering
  \includegraphics[width=0.82\textwidth]{\detokenize{figures/fig-q4-hex37.png}}
  \caption{HEX37共37个覆盖点及从原点出发的800米边Hamilton路，虚线圆为半径1800米的目标域，坐标轴单位为米。}
  \label{fig-q4-hex37}
\end{figure}

\begin{figure}[!htbp]
  \centering
  \includegraphics[width=0.82\textwidth]{\detokenize{figures/fig-q4-paired-tk.png}}
  \caption{离线30组配对平均T/K：SQUARE81为1313.16秒，HEX37为836.91秒，斜线与网格编码覆盖方案。}
  \label{fig-q4-paired-tk}
\end{figure}

演练最新PREFIX\_A记录20260911-231856-q4-p4的T/K为822.38秒，K/N=11/11并完成证书。演练随机非配对，不能与离线36.27\%降幅混为同一对照。

\begin{table}[!htbp]
  \centering
  \caption{问题4演练摘录（非正式）}
  \label{tab-q4-practice}
  \begin{tabularx}{\textwidth}{>{\centering\arraybackslash}X>{\centering\arraybackslash}X>{\centering\arraybackslash}X>{\centering\arraybackslash}X>{\centering\arraybackslash}X}
    \toprule
    证据 & 路线 & K/N & T/K (s) & 证书 \\
    \midrule
    20260911-231856-q4-p4 & PREFIX\_A & 11/11 & 822.38 & 是 \\
20260911-231927-q4-p4 & PREFIX\_A & 10/10 & 930.28 & 是 \\
    \bottomrule
  \end{tabularx}
\end{table}

问题3与问题4正式测试各3次，中止也占次。当前无正式日志，表1不能填写案例编码或成绩。

\begin{table}[!htbp]
  \centering
  \caption{表1 正式测试结果（缺失）}
  \label{tab-formal-missing}
  \begin{tabularx}{\textwidth}{>{\centering\arraybackslash}X>{\centering\arraybackslash}X>{\centering\arraybackslash}X>{\centering\arraybackslash}X>{\centering\arraybackslash}X}
    \toprule
    测试 & 案例编码 & 清除干扰源个数 & 平均定位清除时间 & 程序运行时间 \\
    \midrule
    问题3测试1 & 无正式日志 & 无正式日志 & 无正式日志 & 无正式日志 \\
问题3测试2 & 无正式日志 & 无正式日志 & 无正式日志 & 无正式日志 \\
问题3测试3 & 无正式日志 & 无正式日志 & 无正式日志 & 无正式日志 \\
问题4测试1 & 无正式日志 & 无正式日志 & 无正式日志 & 无正式日志 \\
问题4测试2 & 无正式日志 & 无正式日志 & 无正式日志 & 无正式日志 \\
问题4测试3 & 无正式日志 & 无正式日志 & 无正式日志 & 无正式日志 \\
    \bottomrule
  \end{tabularx}
\end{table}

\section{验证与检验}

问题1用解析反例与代码演示交叉检验：直径定义与最小覆盖圆分离，覆盖判定与几何事实一致。问题2把c1与J(q)限制为候选推荐，避免把演示点写成现场最优。

问题3以官方演练全清率为主检验，GREEDY\_FAST为10/10，平均T/K为341秒。ABORT对照表明去掉SAFE会破坏全清，因此主线保留兜底。各局非配对，区间244至450秒只描述该批演练，不外推总体分布。

问题4离线几何使用与题设一致的有界误差与计时，30/30全清且证书，HEX37平均T/K为836.91秒。演练PREFIX\_A局T/K为822.38秒，证书完成。独立极端误差留出在修订记录中为另外30/30，本文主结论以paired\_summary.json为准。

\section{敏感性与稳健性}

覆盖点集从SQUARE81换到HEX37是一次结构扰动：测向次数均值由679.73降至337.10，路程由68250.75米降至46196.45米，T/K下降36.27\%，全清证书仍为30/30。说明在当前离线误差模型下，证书点数减少没有破坏可检测性。

问题3对SAFE开关敏感：ABORT全清仅1/10，故不能为刷T/K关闭兜底。问题2对网格中心敏感：固定20至3000米极坐标网格已被删除，改以F1最小覆盖圆为局部中心，避免把搜索域写成与第一点示向无关的圆环。

示向度两位小数量化在实现中使用1.005度数值包络，物理误差常数仍为1度，参数表未改。该包络只吸收量化，不作为新的题设误差。

\section{模型局限}

正式测试未执行，表1成绩缺失。演练与离线几何不能代替正式三次测试，也不能把演练T/K写入正式平均定位清除时间。

HEX37未证明是最少点数，也未声称路径全局最优。七点发现集不是最少覆盖。GREEDY\_FAST不含联合路线全局最优，也不含非凸失败域扣除。

Q4的no\_signal不能推出距离大于1000米。光学兜底保证的是当前凸定位域可被20米球覆盖，若观测与模型矛盾则报告不一致，而不是把频道判空。

\section{结论}

交会定位应输出可行集直径与状态，直径圆不能覆盖定位区域。第二检测点应按最坏剩余覆盖半径在第一点可行域附近选取。全向在线主线GREEDY\_FAST在10局演练中全部清除，平均T/K为341秒。混合源采用HEX37证书网格与累计定位，离线30/30全清，演练PREFIX\_A局T/K为822.38秒。

权衡上，保留SAFE以换全清，用更稀的三角晶格换行程，但不牺牲证书。实施上应继续只跑演练直至正式授权；提交前补齐表1与未改名日志。当前结论的边界是演练与离线几何，不是正式测试分布。

\phantomsection
\label{mcm:main-body-end}

\clearpage
\section*{参考文献}
\begin{enumerate}[label={[\arabic*]},leftmargin=*,itemsep=0.5em]
  \item Church R, ReVelle C. The maximal covering location problem. Papers in Regional Science, 1974, 32(1): 101-118.
  \item Dijkstra E W. A note on two problems in connexion with graphs. Numerische Mathematik, 1959, 1(1): 269-271.
  \item Galceran E, Carreras M. A survey on coverage path planning for robotics. Robotics and Autonomous Systems, 2013, 61(12): 1258-1276.
  \item Bishop A N, Fidan B, Anderson B D O, et al. Optimality analysis of sensor-target localization geometries. Automatica, 2010, 46(3): 479-492.
  \item Stansfield R G. Statistical theory of D.F. fixing. Journal of the Institution of Electrical Engineers, 1947, 94(15): 762-770.
\end{enumerate}
\end{document}
